\section{Methodology}
\label{sec:method}
% [6 pts] Explain SDS, VSD, and SDI; both the concept and your implementation.

\noindent\textbf{Background.}
All three methods share the same frozen Stable Diffusion UNet $\epsilon_\theta(x_t, y, t)$~\cite{rombach2022ldm}, which predicts the noise in a noisy latent $x_t$ given a prompt $y$ and timestep $t$.
The forward process~\cite{ho2020ddpm} produces $x_t = \sqrt{\bar\alpha_t}\,x_0 + \sqrt{1-\bar\alpha_t}\,\epsilon$ with $\epsilon\sim\mathcal N(0, I)$.
Noise predictions use classifier-free guidance (CFG)~\cite{ho2022cfg}, $\hat\epsilon_\theta = \epsilon_{\text{uncond}} + s\,(\epsilon_{\text{cond}} - \epsilon_{\text{uncond}})$, where the unconditional branch uses the negative prompt ``low quality''.
The optimized parameter is the latent $x_0$, initialized from $\mathcal N(0, I)$ and decoded with the VAE only for visualization.
The optimizer (AdamW, cosine schedule, gradient clipping) is the one provided in \texttt{main.py}.

%-------------------------------------------------------------------------
\subsection{Score Distillation Sampling (SDS)}
\label{sec:sds}

\noindent\textbf{Concept.}
SDS~\cite{poole2023dreamfusion} noises the current latent with a random $t$ and $\epsilon$, asks the diffusion model to predict the noise, and uses the difference as the update direction:
$\nabla_{x_0} L_{\text{SDS}} = w(t)\,(\hat\epsilon_\theta(x_t, y, t) - \epsilon)$.
Intuitively, the residual $\hat\epsilon_\theta - \epsilon$ points from the current image towards what the model thinks a clean image for this prompt looks like.
This is the gradient of the ordinary diffusion loss with the UNet Jacobian term dropped; skipping it makes SDS cheap (no backpropagation through the UNet) and, as \cref{sec:diff_loss} shows, is also what makes it work at all.

\noindent\textbf{Implementation.}
I sample $t\sim\mathcal U[20, 980]$ and $\epsilon$, add noise with the scheduler, and predict the noise with CFG under \texttt{no\_grad}.
The gradient uses the DreamFusion weight $w(t)=1-\bar\alpha_t$.
To feed this hand-made gradient to the optimizer, I write it as an MSE to a detached target $x_0 - \text{grad}$, whose derivative with respect to $x_0$ is exactly the SDS gradient.

\begin{lstlisting}
def get_sds_loss(self, latents, text_embeddings, guidance_scale=7.5):
    B = latents.shape[0]
    t = torch.randint(self.min_step, self.max_step + 1, (B,), device=self.device)
    noise = torch.randn_like(latents)
    with torch.no_grad():
        latents_noisy = self.scheduler.add_noise(latents, noise, t)
        noise_pred = self.get_noise_preds(latents_noisy, t, text_embeddings, guidance_scale)
    w = (1 - self.alphas[t]).view(B, 1, 1, 1)
    grad = torch.nan_to_num(w * (noise_pred - noise))
    # d(loss)/d(latents) == grad (skips the UNet Jacobian)
    target = (latents - grad).detach()
    return 0.5 * F.mse_loss(latents, target, reduction="sum") / B
\end{lstlisting}

%-------------------------------------------------------------------------
\subsection{Variational Score Distillation (VSD)}
\label{sec:vsd}

\noindent\textbf{Concept.}
In SDS the reference $\epsilon$ is pure random noise, so the update is noisy and SDS needs a very large guidance scale to converge, which causes over-smoothed and over-saturated results.
VSD~\cite{wang2023prolificdreamer} replaces $\epsilon$ with the prediction of a second network $\epsilon_\phi$, a LoRA~\cite{hu2022lora} on top of the same UNet, which is trained online to model the distribution of the images currently being generated.
The update $w(t)\,(\hat\epsilon_\theta - \epsilon_\phi)$ then compares ``what the pretrained model expects'' with ``what we currently generate'', which is a much cleaner signal and works with a normal guidance scale of 7.5.

\noindent\textbf{Implementation.}
The LoRA (rank 4 on the attention projections) is set up by the provided code.
For $\epsilon_\theta$ I temporarily disable the adapters and predict with CFG; for $\epsilon_\phi$ I enable them and use only the conditional branch.
Both predictions are made without gradients, and the VSD term uses the same detached-target trick as SDS, so it only updates the latent.
The LoRA is trained with the standard denoising loss on a \emph{detached} copy of the latent with a fresh $t$ and $\epsilon$, so it only updates the LoRA.
The returned loss is the sum of both terms, and \texttt{main.py} steps the latent and the LoRA parameters together.

\begin{lstlisting}
def get_vsd_loss(self, latents, text_embeddings, guidance_scale=7.5, lora_loss_weight=1.0):
    B = latents.shape[0]
    cond = text_embeddings.chunk(2)[1]
    t = torch.randint(self.min_step, self.max_step + 1, (B,), device=self.device)
    noise = torch.randn_like(latents)
    with torch.no_grad():
        latents_noisy = self.scheduler.add_noise(latents, noise, t)
        self.unet.disable_adapters()           # eps_theta: pretrained UNet + CFG
        eps_pre = self.get_noise_preds(latents_noisy, t, text_embeddings, guidance_scale)
        self.unet.enable_adapters()            # eps_phi: LoRA UNet, conditional only
        eps_lora = self.unet(latents_noisy, t, encoder_hidden_states=cond).sample
    w = (1 - self.alphas[t]).view(B, 1, 1, 1)
    grad = torch.nan_to_num(w * (eps_pre - eps_lora))
    target = (latents - grad).detach()
    loss_vsd = 0.5 * F.mse_loss(latents, target, reduction="sum") / B

    # LoRA learns the distribution of the current latents
    t2 = torch.randint(self.min_step, self.max_step + 1, (B,), device=self.device)
    noise2 = torch.randn_like(latents)
    noisy2 = self.scheduler.add_noise(latents.detach(), noise2, t2)
    pred2 = self.unet(noisy2, t2, encoder_hidden_states=cond).sample
    loss_lora = F.mse_loss(pred2.float(), noise2.float())
    return loss_vsd + lora_loss_weight * loss_lora
\end{lstlisting}

%-------------------------------------------------------------------------
\subsection{Score Distillation via Inversion (SDI)}
\label{sec:sdi}

\noindent\textbf{Concept.}
SDI~\cite{lukoianov2024sdi} keeps the SDS framework but fixes its noise: instead of a fresh random $\epsilon$ at every step, the noisy latent $x_t$ is obtained by running DDIM~\cite{song2021ddim} \emph{backwards} from the current $x_0$.
The noise is therefore consistent with the current image, so denoising $x_t$ gives a target $\hat x_0$ that is a sharpened version of the current latent rather than an average over many random noises.
The latent is then pulled towards $\hat x_0$ with a simple MSE.
SDI also replaces the random timestep of SDS with a schedule that goes from large to small $t$, so the image first settles its layout and then refines its details.

\noindent\textbf{Implementation.}
\texttt{invert\_noise()} is a hand-written deterministic DDIM inversion ($\eta = 0$, as allowed): it splits $[0, t]$ into $n$ steps and, at each step, predicts the noise with CFG scale $-7.5$, estimates $\hat x_0$, and re-noises it to the next (larger) timestep.
In \texttt{get\_sdi\_loss()}, $t$ is annealed linearly from 980 to 20 over the 1000 iterations.
Every \texttt{update\_interval} iterations the current latent is inverted to $x_t$, the noise is predicted with CFG 7.5, and a one-step denoise gives the target, which is cached.
In between, the loss is only the MSE to the cached target, so most iterations need no UNet call at all.

\begin{lstlisting}
@torch.no_grad()
def invert_noise(self, latents, target_t, text_embeddings, guidance_scale=-7.5, n_steps=10, eta=0.3):
    B = latents.shape[0]
    target = int(target_t.max().item())
    timesteps = torch.linspace(0, target, n_steps + 1, device=self.device).round().long()
    x = latents.clone()
    for t_cur, t_next in zip(timesteps[:-1], timesteps[1:]):
        a_cur, a_next = self.alphas[t_cur], self.alphas[t_next]
        eps = self.get_noise_preds(x, t_cur.expand(B), text_embeddings, guidance_scale)
        x0 = (x - (1 - a_cur).sqrt() * eps) / a_cur.sqrt()
        x = a_next.sqrt() * x0 + (1 - a_next).sqrt() * eps
    return x
\end{lstlisting}

\begin{lstlisting}
def get_sdi_loss(self, latents, text_embeddings, guidance_scale=7.5,
                 current_iter=0, total_iters=500, inversion_guidance_scale=-7.5,
                 inversion_n_steps=10, inversion_eta=0.3, update_interval=25):
    B = latents.shape[0]
    progress = current_iter / max(1, total_iters)        # anneal t: 980 -> 20
    t_val = int(round(self.max_step - (self.max_step - self.min_step) * progress))
    t = torch.full((B,), t_val, dtype=torch.long, device=self.device)

    if current_iter % update_interval == 0 or not hasattr(self, 'sdi_target'):
        with torch.no_grad():
            x_t = self.invert_noise(latents, t, text_embeddings, inversion_guidance_scale,
                                    inversion_n_steps, inversion_eta)
            eps = self.get_noise_preds(x_t, t, text_embeddings, guidance_scale)
            a_t = self.alphas[t].view(B, 1, 1, 1)
            target = (x_t - (1 - a_t).sqrt() * eps) / a_t.sqrt()   # one-step denoise
            self.sdi_target = torch.nan_to_num(target).detach()
    return 0.5 * F.mse_loss(latents, self.sdi_target, reduction="sum") / B
\end{lstlisting}
